\begin{figure*}[t]
\centering
\tikzset{
  op/.style={draw, rounded corners, align=center, font=\scriptsize\sffamily,
             minimum width=2.1cm, minimum height=0.5cm, fill=white},
  top/.style={draw, rounded corners, align=center, font=\scriptsize\sffamily,
              minimum width=4.6cm, fill=black!5},
  flow/.style={-{Latex[length=1.5mm]}},
  blabel/.style={font=\scriptsize\bfseries\sffamily, inner sep=1pt}}

% ---- (a) distance-predicate join ----
\begin{minipage}[t]{0.48\textwidth}
\centering
\begin{lstlisting}[language=SQL, basicstyle=\ttfamily\scriptsize, frame=none,
                   aboveskip=0pt, belowskip=4pt]
SELECT *, array_distance(l.v, r.v) AS dist
FROM l JOIN r ON array_distance(l.v, r.v) <= 5;
\end{lstlisting}
\begin{tikzpicture}
  \node[op] (s0) at (0,0)     {GPU\_SCAN};
  \node[op] (s3) at (2.7,0)   {GPU\_SCAN};
  \node[op] (p1) at (0,1.3)   {PARTITION};
  \node[op] (p4) at (2.7,1.3) {PARTITION};
  \node[op] (c2) at (0,2.6)   {CONCAT};
  \node[op] (c5) at (2.7,2.6) {CONCAT};
  \node[top] (t) at (1.35,4.3)
    {RESULT\_COLLECTOR\\ PROJECTION\\ \textbf{VECTOR\_THRESHOLD\_JOIN}\\[1pt]
     {\scriptsize type: INNER}};
  \draw[flow] (s0) -- (p1);
  \draw[flow] (s3) -- (p4);
  \draw[flow] (p1) -- (c2);
  \draw[flow] (p4) -- (c5);
  \draw[flow] (c2) -- (t.south -| c2) node[blabel, midway, left]  {FULL};
  \draw[flow] (c5) -- (t.south -| c5) node[blabel, midway, right] {FULL};
\end{tikzpicture}

{\small (a) distance-predicate join}
\end{minipage}
\hfill
% ---- (b) equi-join ----
\begin{minipage}[t]{0.48\textwidth}
\centering
\begin{lstlisting}[language=SQL, basicstyle=\ttfamily\scriptsize, frame=none,
                   aboveskip=0pt, belowskip=4pt]
SELECT l.id, r.id
FROM l JOIN r ON l.id = r.id;
\end{lstlisting}
\begin{tikzpicture}
  \node[op] (s0) at (0,0)     {GPU\_SCAN};
  \node[op] (s3) at (2.7,0)   {GPU\_SCAN};
  \node[op] (p1) at (0,1.3)   {PARTITION};
  \node[op] (p4) at (2.7,1.3) {PARTITION};
  \node[op] (c2) at (0,2.6)   {CONCAT};
  \node[op] (c5) at (2.7,2.6) {CONCAT};
  \node[top] (t) at (1.35,4.3)
    {RESULT\_COLLECTOR\\ PROJECTION\\ \textbf{HASH\_JOIN}\\[1pt]
     {\scriptsize type: INNER}};
  \draw[flow] (s0) -- (p1);
  \draw[flow] (s3) -- (p4);
  \draw[flow] (p1) -- (c2);
  \draw[flow] (p4) -- (c5);
  \draw[flow] (c2) -- (t.south -| c2) node[blabel, midway, left]  {PARTIAL};
  \draw[flow] (c5) -- (t.south -| c5) node[blabel, midway, right] {PARTIAL};
\end{tikzpicture}

{\small (b) equi-join}
\end{minipage}
\caption{The same query shape under two join predicates. The distance-predicate
join~(a) and the equi-join~(b) produce identical operator trees; only the join
node and its input barriers differ. The vector threshold join consumes both
sides in full because it walks every batch pair, whereas the hash join streams
its probe side with a partial barrier.}
\label{fig:native-plans}
\end{figure*}
